This paper investigates the cognitive demands imposed by technical interview formats, particularly focusing on whiteboard-based problem solving commonly used in software engineering hiring. It highlights how such interview settings may unintentionally increase cognitive load due to factors like lack of supportive tools, public performance pressure, and enforced verbalization of thought processes. These factors can interfere with a candidate's ability to effectively reason and solve problems, potentially disadvantaging otherwise capable individuals. The study is motivated by concerns that traditional interview practices may not accurately reflect a candidate's true abilities and could contribute to biased hiring outcomes.

To explore this, the authors propose a novel methodology using head-mounted eye-tracking technology to measure cognitive load during problem-solving tasks in different interview environments. Through a pilot study comparing whiteboard and paper-based settings, they demonstrate that whiteboard interviews are associated with higher cognitive load, reflected in shorter attention spans, increased regressions, and elevated stress indicators. These findings suggest that interview format significantly impacts candidate performance and cognitive state, emphasizing the need for more inclusive and cognitively considerate interview practices.

\section{Motivation}
The motivation for this work stems from the widespread use of technical interviews—particularly whiteboard based programming tasks—as a primary mechanism for evaluating software engineering candidates. These interviews are intended to provide visibility into a candidate's cognitive processes and problem-solving abilities \cite{8_Ford2017}. However, despite their prevalence, there is growing concern that such formats impose additional cognitive demands that are unrelated to the candidate's actual programming competence. Specifically, the total cognitive load experienced during interviews may be influenced not only by the task itself but also by external factors such as the medium of interaction and the interview setting \cite{14_Paas_1994}. This raises questions about whether current interview practices truly measure the intended skills.

A key motivating factor is the recognition that whiteboard interviews introduce several sources of extraneous cognitive load. The lack of familiar development tools and affordances—such as syntax highlighting—can hinder a candidate’s ability to effectively structure and debug their thoughts \cite{24_Waestlund2005}. Additionally, the interview environment often involves public performance, time pressure, and self-evaluation, all of which contribute to increased stress levels \cite{11_Klitzman1990WorkSN}. Candidates are also typically required to verbalize their thought processes while solving problems, which has been shown to further increase cognitive load and disrupt natural reasoning patterns \cite{15_Parnin_2010,19_Russo_1989}. These factors collectively create an artificial problem-solving environment that may not align with how developers typically work.

Another important motivation is the potential for these cognitive mismatches to introduce bias and inequity into hiring processes. When interview formats are not aligned with diverse cognitive styles, they may disproportionately disadvantage certain groups of candidates \cite{13_Morris_2015}. Prior work has shown that software tools and environments can favor specific problem-solving approaches, potentially reinforcing systemic disparities \cite{3_Burnett2016}. Furthermore, evaluation criteria such as confidence and visible composure during interviews may lead to biased decisions, where less stressed candidates are favored even when performance is equivalent \cite{8_Ford2017}. As a result, traditional technical interviews risk overlooking skilled candidates who may struggle under these conditions, motivating the need to better understand and redesign interview practices.

\section{Methodology}
\subsection{Eye Tracking}
Eye-tracking is the process of locating the eye position and measuring the eye movement of a subject. Eye-trackers are designed to monitor and collect eye-movement data while the subject is looking at a stimulus during an experiment. Head-mounted eye-trackers have embedded infrared cameras and are worn by subjects. In contrast, remote eye-trackers consist of a screen with sensors, where subjects sit in front of the device, such as those used to study eye movements of developers while using IDEs.
\paragraph{Common Eye Tracking Terms} Some common eye-tracking terms include:
\begin{itemize}
    \item \textbf{Areas of Interest (AOI)}: A fixed region corresponding to an object of study.
    \item \textbf{Fixation or Visual Intake (VI)}: Stabilization of the eye on a particular point of the stimulus, typically between 200 to 300 ms. Fixations can be acquired from the eye-tracker in the form of a timestamp and x-y pixel coordinates. Privitera et al. suggest that most of the cognitive processes and information acquisition happen during fixations \cite{16_Privitera2000}.
    \item \textbf{Saccade}: A sudden, rapid eye movement which occurs between two fixations, typically lasting between 40–50 ms. Prior studies claim that cognitive processing and information acquisition are not notable during saccades \cite{6_Duchowski2003,16_Privitera2000}.
    \item \textbf{Regression}: A backward movement of the eye from an AOI to a previously visited AOI.
    \item \textbf{Cognitive Load}: Several studies report eye-tracking measures that reveal cognitive load. Chen et al. conclude that fixation duration and fixation rate are indicators of an increase in attention \cite{4_Chen2011}. The increase or decrease of fixation duration depends on the characteristics of the task. Fixation duration will decrease in stressful tasks that require rapid responses (such as driving a car \cite{23_Falkmer2005} or airplane piloting). In contrast, in tasks that require more cognitive processing, such as reading texts of increasing complexity, fixation duration will increase.
\end{itemize}

\subsection{Approach}
\subsubsection{Placing Boundaries for the Coding Area}
Since the tracker produces pixel-based gaze coordinates, ArUco markers are placed in the environment to map these coordinates to physical locations. These markers are robust to occlusion and angle variations, enabling accurate tracking even with participant movement. The implementation leverages OpenCV to detect markers and associate gaze points with regions in the coding space, forming the foundation for subsequent cognitive load analysis \cite{9_Garrido-Jurado2014}.

\subsubsection{Determining AOIs for Coding Area}
The authors define Areas of Interest (AOIs) by arranging ArUco markers in structured grid patterns on both the whiteboard and paper setups. The grid ensures consistent spatial segmentation, allowing precise mapping of gaze behavior across different regions. Design considerations include maintaining marker visibility regardless of participant distance and enabling systematic computation of gaze transitions. The configuration differs slightly between whiteboard and paper to account for scale and viewing distance, but both aim to support reliable analysis of visual attention patterns.

\subsubsection{Calculating Gaze Regressions}
For the whiteboard setting, regressions are defined as backward gaze movements across grid rows, while for the paper setting, pixel distances are converted into physical measurements to account for variable viewing distances. A threshold-based approach is used to identify meaningful regressions, enabling comparison across conditions. These regressions serve as a proxy for memory difficulty and uncertainty during problem solving.

\subsubsection{Pilot Methodology}
The pilot study involves 11 participants who complete programming tasks under both whiteboard and paper conditions while wearing eye-tracking devices. Tasks are selected from a standard interview preparation resource and are comparable in difficulty. A Latin Square design is used to control for ordering effects, and participants are allowed to use any programming language. The study collects multiple eye-tracking metrics alongside qualitative observations, providing initial evidence on how different interview settings influence cognitive load and stress.

\section{Preliminary Results}
\subsection{Under pressure}
The authors observe significantly shorter fixation durations when participants work on the whiteboard, suggesting that candidates are forced to process information more quickly and are unable to sustain attention for long periods. This reduced attention span limits opportunities for deeper reflection and careful reasoning. In contrast, participants working on paper tend to take slightly longer and appear more composed. Qualitative observations reinforce these findings, as participants reported higher stress levels in the whiteboard condition and exhibited visible nervous behaviors such as fidgeting and involuntary movements. While a moderate level of stress may occasionally enhance performance, the study suggests that excessive pressure in whiteboard interviews can hinder cognitive effectiveness and negatively impact problem-solving quality.

\subsection{Higher cognitive load}
Participants exhibited significantly more gaze regressions when using the whiteboard, indicating frequent backtracking and difficulty maintaining information in working memory. This suggests increased uncertainty and cognitive strain during the task. Additional eye-tracking metrics, such as higher saccade duration and velocity, point to elevated stress levels and reduced concentration. The study also finds that blink duration increases substantially in the whiteboard condition, which is associated with more demanding cognitive processing. Some of these effects are attributed to the inherent stress of public performance, while others may stem from the larger and less structured workspace of the whiteboard, which makes it harder to locate and manage code visually. Overall, the results demonstrate that whiteboard settings create a more cognitively taxing environment that can impair performance.
